% Unit 040 English translation of chapter3.tex lines 1587--1709.
% Source: Methods of Algebra, Volume 2, commit
% 9a5803ff77dd3257484cb177f851a73770a59dd3 (CC BY 4.0).
% stable-unit-id: o014.aljabr2.chapter3.truncation-functors
% Coverage: Section 3.9 in full.

% segment-id: o014.li2.chapter3.truncation-functors.h001
\section{Truncation Functors}\label{sec:truncation-functors}
\hypertarget{o014.aljabr2.chapter3.truncation-functors}{ }

% segment-id: o014.li2.chapter3.truncation-functors.p001
At the beginning of this section, let $\mathcal{A}$ be an additive category.

% segment-id: o014.li2.chapter3.truncation-functors.de001
\begin{definition}\label{def:cplx-cat-variant}
	\index{complex!bounded above, bounded below, bounded}
	\index[sym1]{CbA@$\cate{C}^+(\mathcal{A})$, $\cate{C}^-(\mathcal{A})$, $\cate{C}^{\bdd}(\mathcal{A})$}
	\index[sym1]{CstA@$\cate{C}^{[s,t]}(\mathcal{A})$, $\cate{C}^{\geq s}(\mathcal{A})$, $\cate{C}^{\leq t}(\mathcal{A})$}
	For a complex $X \in \Obj(\cate{C}(\mathcal{A}))$, use the following
	terms and corresponding full subcategories:
% segment-id: o014.li2.chapter3.truncation-functors.q001
	\begin{center}\small\setlength{\tabcolsep}{3pt}\begin{tabular}{|c|c|c|c|}\hline
		Term & bounded & bounded below & bounded above \\
		Condition & $|n| \gg 0 \implies X^n = 0 $ & $n \ll 0 \implies X^n = 0$ & $n \gg 0 \implies X^n = 0$ \\
		Full subcategory & $\cate{C}^{\bdd}(\mathcal{A})$ & $\cate{C}^+(\mathcal{A})$ & $\cate{C}^-(\mathcal{A})$ \\ \hline
	\end{tabular}\end{center}

	More generally, for $-\infty \leq s \leq t \leq +\infty$, let
	$\cate{C}^{[s,t]}(\mathcal{A})$ denote the full subcategory consisting of
	complexes satisfying $n \notin [s, t] \implies X^n = 0$, and define
% segment-id: o014.li2.chapter3.truncation-functors.q002
	\[ \cate{C}^{\geq s}(\mathcal{A}) := \cate{C}^{[s, +\infty]}(\mathcal{A}), \quad \cate{C}^{\leq t}(\mathcal{A}) := \cate{C}^{[-\infty, t]}(\mathcal{A}). \]
\end{definition}

% segment-id: o014.li2.chapter3.truncation-functors.p002
All these full subcategories are additive categories. If $\mathcal{A}$ is an
abelian category, then $\cate{C}^+(\mathcal{A})$,
$\cate{C}^-(\mathcal{A})$, and
$\cate{C}^{\bdd}(\mathcal{A}) = \cate{C}^+(\mathcal{A}) \cap
\cate{C}^-(\mathcal{A})$ are abelian subcategories.\footnote{Translator's
note (O014-C050): the source states the second assertion without requiring
$\mathcal{A}$ to be abelian, although at this point $\mathcal{A}$ has only
been assumed additive. The assertion is made conditional here.} The shift
functor $[n]$ preserves $\cate{C}^+(\mathcal{A})$,
$\cate{C}^-(\mathcal{A})$, and $\cate{C}^{\bdd}(\mathcal{A})$, but maps
$\cate{C}^{[a,b]}(\mathcal{A})$ to
$\cate{C}^{[a-n, b-n]}(\mathcal{A})$.

% segment-id: o014.li2.chapter3.truncation-functors.p003
Moreover, $\mathcal{A}$ is naturally identified with
$\cate{C}^{\geq 0}(\mathcal{A}) \cap \cate{C}^{\leq 0}(\mathcal{A})$.

% segment-id: o014.li2.chapter3.truncation-functors.p004
For every $\star \in \{+, -, \bdd\}$, the notion of homotopy between
morphisms (Definition~\ref{def:homotopy}) restricts to
$\cate{C}^{\star}(\mathcal{A})$. Thus
Definition~\ref{def:cplx-homotopy-cat} has the following variant.

% segment-id: o014.li2.chapter3.truncation-functors.de002
\begin{definition}\label{def:cplx-homotopy-cat-variant}
	\index[sym1]{KbA@$\cate{K}^+(\mathcal{A})$, $\cate{K}^-(\mathcal{A})$, $\cate{K}^{\bdd}(\mathcal{A})$}
	For $\star \in \{+, -, \bdd\}$, $\cate{K}(\mathcal{A})$ has a full
	additive subcategory $\cate{K}^\star(\mathcal{A})$ that is closed under
	the shift functors.
\end{definition}

% segment-id: o014.li2.chapter3.truncation-functors.p005
From now on, we require $\mathcal{A}$ to be an abelian category. We study how
to truncate the part of a complex $X$ in degrees $<n$ (or $>n$). The naive
idea is to replace every term in that part by $0$. Although simple, this
disrupts cohomology, so we introduce a more refined version.

% segment-id: o014.li2.chapter3.truncation-functors.de003
\begin{definition}[Truncation functors]\label{def:truncation-cplx}
	\index{truncation functor}
	\index[sym1]{taun@$\tau^{\leq n}$, $\tau^{\geq n}$}
	\index[sym1]{tautilden@$\tilde{\tau}^{\leq n}$, $\tilde{\tau}^{\geq n}$}
	Let $\mathcal{A}$ be an abelian category and let $n \in \Z$. For a
	complex $X$, define
% segment-id: o014.li2.chapter3.truncation-functors.g001
	\[\begin{tikzcd}
		[column sep=small, every cell/.append style = {font = \small}]
		\tau^{\leq n} X \arrow[phantom, r, "" {name=A, yshift=2ex}] \arrow[hookrightarrow, d] & \cdots \arrow[r] & X^{n-2} \arrow[r] & X^{n-1} \arrow[r] & \Ker(d^n) \arrow[r] \arrow[hookrightarrow, d] & 0 \arrow[r] & 0 \arrow[r] & \cdots \\
		\tilde{\tau}^{\leq n} X \arrow[hookrightarrow, d] & \cdots \arrow[r] & X^{n-2} \arrow[r] & X^{n-1} \arrow[r] & X^n \arrow[r] & \Coim(d^n) \arrow[r] \arrow[hookrightarrow, d] & 0 \arrow[r] & \cdots \\
		X \arrow[twoheadrightarrow, d] & \cdots \arrow[r] & X^{n-2} \arrow[r] & X^{n-1} \arrow[r] \arrow[twoheadrightarrow, d] & X^n \arrow[r] & X^{n+1} \arrow[r] & X^{n+2} \arrow[r] & \cdots \\
		\tilde{\tau}^{\geq n} X \arrow[twoheadrightarrow, d] & \cdots \arrow[r] & 0 \arrow[r] & \Image(d^{n-1}) \arrow[r] & X^n \arrow[r] \arrow[twoheadrightarrow, d] & X^{n+1} \arrow[r] & X^{n+2} \arrow[r] & \cdots \\
		\tau^{\geq n} X \arrow[phantom, r, "" {name=B, yshift=-2ex}] & \cdots \arrow[r] & 0 \arrow[r] & 0 \arrow[r] & \Coker(d^{n-1}) \arrow[r] & X^{n+1} \arrow[r] & X^{n+2} \arrow[r] & \cdots
		\arrow[dash, thick, from=A, to=B]
	\end{tikzcd}\]
	The omitted terms are clear, and in the vertical direction only morphisms
	other than $\identity$ and $0$ are marked. We thus obtain morphisms between
	the complexes listed on the left; in this notation, every $\Coim(d^i)$ is
	identified with $\Image(d^i)$ through the canonical comparison
	isomorphism.
\end{definition}

% segment-id: o014.li2.chapter3.truncation-functors.p006
It is clear that $\tau^{\leq n}$, $\tilde{\tau}^{\leq n}$,
$\tau^{\geq n}$, and $\tilde{\tau}^{\geq n}$ are all additive functors.
We have $\tau^{\leq n} = [-n] \circ \tau^{\leq 0} \circ [n]$; the same
holds for $\tilde{\tau}^{\leq n}$, $\tau^{\geq n}$, and
$\tilde{\tau}^{\geq n}$.

% segment-id: o014.li2.chapter3.truncation-functors.p007
Moreover, if $m \leq n$, there are evident epimorphisms
$\tau^{\geq m} X \twoheadrightarrow \tau^{\geq n} X$ and
$\tilde{\tau}^{\geq m} X \twoheadrightarrow \tilde{\tau}^{\geq n} X$,
and evident monomorphisms
$\tau^{\leq m} X \hookrightarrow \tau^{\leq n} X$ and
$\tilde{\tau}^{\leq m} X \hookrightarrow \tilde{\tau}^{\leq n} X$.

% segment-id: o014.li2.chapter3.truncation-functors.le001
\begin{lemma}\label{prop:truncation-ses}
	For every $k \in \Z$, the comparison morphisms
	$\tau^{\leq n}X \to \tilde{\tau}^{\leq n}X$ and
	$\tilde{\tau}^{\geq n}X \to \tau^{\geq n}X$ in the diagram above
	induce the following isomorphisms in $\mathcal{A}$\footnote{Translator's
	note (O014-C051): the source refers to ``the morphisms above,'' which could
	be read as the transition morphisms for $m\leq n$ in the immediately
	preceding paragraph; only the two comparison morphisms displayed here
	induce these cohomology isomorphisms.}
% segment-id: o014.li2.chapter3.truncation-functors.q003
	\begin{align*}
		\Hm^k\left( \tau^{\leq n} X\right) & \rightiso \Hm^k\left( \tilde{\tau}^{\leq n} X\right) \simeq \begin{cases} \Hm^k(X), & k \leq n \\ 0, & k > n \end{cases}, \\
		\Hm^k\left( \tilde{\tau}^{\geq n} X\right) & \rightiso \Hm^k\left( \tau^{\geq n} X\right) \simeq \begin{cases} \Hm^k(X), & k \geq n \\ 0, & k < n \end{cases},
	\end{align*}
	and the following short exact sequences in $\cate{C}(\mathcal{A})$:
% segment-id: o014.li2.chapter3.truncation-functors.q004
	\begin{equation*}\begin{gathered}
		0 \to \tilde{\tau}^{\leq n-1} X \to \tau^{\leq n} X \to \Hm^n(X)[-n] \to 0, \\
		0 \to \Hm^n(X)[-n] \to \tau^{\geq n} X \to \tilde{\tau}^{\geq n+1} X \to 0, \\
		0 \to \tau^{\leq n} X \to X \to \tilde{\tau}^{\geq n+1} X \to 0, \\
		0 \to \tilde{\tau}^{\leq n-1} X \to X \to \tau^{\geq n} X \to 0, \\
		0 \to \tau^{\leq n} X \to \tilde{\tau}^{\leq n} X \to \Cone\left(\identity_{\Image(d_X^n) [-n-1]} \right) \to 0.
	\end{gathered}\end{equation*}
	Here $\Hm^n(X)[-n]$ is understood according to
	Convention~\ref{con:concentrated-cplx}. All these morphisms are functorial
	in $X$.
\end{lemma}
% segment-id: o014.li2.chapter3.truncation-functors.pf001
\begin{proof}
	The degree-$n$ component of the morphism
	$\tilde{\tau}^{\leq n-1} X \to \tau^{\leq n} X$ is
	$\Coim(d^{n-1}) \xrightarrow{\sim}
	\Image(d^{n-1}) \hookrightarrow \Ker(d^n)$,\footnote{Translator's note
	(O014-C052): the source writes $\Image(d^{n-1})$ directly as the
	degree-$n$ term of $\tilde{\tau}^{\leq n-1}X$; that term is actually
	$\Coim(d^{n-1})$, canonically identified with its image.}
	and every other component is $\identity$. The degree-$n$ component of the
	morphism $\tau^{\leq n} X \to \Hm^n(X)[-n]$ is
% segment-id: o014.li2.chapter3.truncation-functors.q005
	\[ \Ker(d^n) \twoheadrightarrow \Ker(d^n)/\Image(d^{n-1}) = \Hm^n(X). \]

	The degree-$n$ component of the morphism
	$\tau^{\geq n} X \to \tilde{\tau}^{\geq n+1} X$ is the epimorphism
	induced by $d^n$, $\Coker(d^{n-1}) \twoheadrightarrow \Image(d^n)$, and
	every other component is $\identity$. The degree-$n$ component of the
	morphism $\Hm^n(X)[-n] \to \tau^{\geq n} X$ is
% segment-id: o014.li2.chapter3.truncation-functors.q006
	\[ \Ker(d^n)/\Image(d^{n-1}) \hookrightarrow X^n/\Image(d^{n-1}) = \Coker(d^{n-1}). \]

	All remaining verifications are routine and are left to the reader as
	exercises.
\end{proof}

% segment-id: o014.li2.chapter3.truncation-functors.p008
Thus $\tau^{\leq n}$ and $\tilde{\tau}^{\leq n}$ (or $\tau^{\geq n}$ and
$\tilde{\tau}^{\geq n}$) do indeed truncate cohomology in degrees $>n$
(or $<n$).

% segment-id: o014.li2.chapter3.truncation-functors.re001
\begin{remark}\label{rem:truncation-sigma}
	For a complex $X \in \Obj(\cate{C}(\mathcal{A}))$, applying the functor
	$\sigma$ of Definition--Proposition~\ref{def:sigma} gives
	\[
	\sigma X \in \Obj(\cate{C}(\mathcal{A}^{\opp})) =
	\Obj(\cate{C}(\mathcal{A}^{\opp})^{\opp}).
	\]
	Ignoring some $\pm$ signs that have no effect in this context, this
	operation amounts to reversing the arrows in the complex and then
	reversing the degree indices. By the duality between $\Ker$ and $\Coker$,
	under this operation $\tau^{\geq n}(X)$ corresponds to
	$\tau^{\leq -n}(\sigma X) \in
	\Obj(\cate{C}(\mathcal{A}^{\opp}))$, and so forth. Similarly, because
	$\Image$ and $\Coim$ are dual to each other,
	$\tilde{\tau}^{\geq n}(X) \in \Obj(\cate{C}(\mathcal{A}))$ corresponds
	to $\tilde{\tau}^{\leq -n}(\sigma X) \in
	\Obj(\cate{C}(\mathcal{A}^{\opp}))$.
\end{remark}

% segment-id: o014.li2.chapter3.truncation-functors.p009
Compared with $\tilde{\tau}^{\leq n}$ and $\tilde{\tau}^{\geq n}$, the
functors $\tau^{\leq n}$ and $\tau^{\geq n}$ have some properties that are
easier to use. The first is adjunction.

% segment-id: o014.li2.chapter3.truncation-functors.pr001
\begin{proposition}\label{prop:truncation-adjunction}
	For every $n \in \Z$, there are the following adjunctions:
% segment-id: o014.li2.chapter3.truncation-functors.q007
	\begin{gather*}
		\Hom_{\cate{C}(\mathcal{A})}(X, Y) \simeq \Hom_{\cate{C}^{\leq n}(\mathcal{A})}\left(X, \tau^{\leq n} Y \right), \quad X \in \Obj(\cate{C}^{\leq n}(\mathcal{A})) \\ \Hom_{\cate{C}(\mathcal{A})}(X, Y) \simeq \Hom_{\cate{C}^{\geq n}(\mathcal{A})}\left(\tau^{\geq n} X, Y \right), \quad Y \in \Obj(\cate{C}^{\geq n}(\mathcal{A})).
	\end{gather*}
\end{proposition}
% segment-id: o014.li2.chapter3.truncation-functors.pf002
\begin{proof}
	This is clear from the definition of truncation. The reader may write down
	the corresponding unit and counit morphisms: they are the morphisms shown
	in Definition~\ref{def:truncation-cplx}, or $\identity$.
\end{proof}

% segment-id: o014.li2.chapter3.truncation-functors.p010
Next, $\tau^{\leq n}$ and $\tau^{\geq n}$ descend to
$\cate{K}(\mathcal{A})$.

% segment-id: o014.li2.chapter3.truncation-functors.dp001
\begin{definition-proposition}\label{def:K-truncation-functors}
	\index[sym1]{taun}
	For every $n \in \Z$, the functor $\tau^{\leq n}$ (or
	$\tau^{\geq n}$) naturally induces a functor from
	$\cate{K}(\mathcal{A})$ to itself, still denoted by $\tau^{\leq n}$ (or
	$\tau^{\geq n}$).
\end{definition-proposition}
% segment-id: o014.li2.chapter3.truncation-functors.pf003
\begin{proof}
	Let $h \in \Hom^{-1}(X, Y)$. For the case of $\tau^{\leq n}$, define
	$\overline{h}^n := h^n \circ \bigl(\Ker(d_X^n) \hookrightarrow X^n\bigr)$.
	Define $\overline{h} \in
	\Hom^{-1}\left( \tau^{\leq n} X, \tau^{\leq n} Y \right)$ by the
	following diagram:
% segment-id: o014.li2.chapter3.truncation-functors.g002
	\[\begin{tikzcd}
		\cdots \arrow[r] & X^{n-2} \arrow[r] \arrow[ld, "h^{n-2}" description] & X^{n-1} \arrow[r] \arrow[ld, "{h^{n-1}}" description] & \Ker(d_X^n) \arrow[r] \arrow[ld, "{\overline{h}^n}" description] & 0 \arrow[r] \arrow[ld] & 0 \arrow[r] \arrow[ld] & \cdots \\
		\cdots \arrow[r] & Y^{n-2} \arrow[r] & Y^{n-1} \arrow[r] & \Ker(d_Y^n) \arrow[r] & 0 \arrow[r] & 0 \arrow[r] & \cdots ;
	\end{tikzcd}\]
	It is easy to see that
	$\tau^{\leq n}\left( d^{-1} h \right) = d^{-1} \overline{h}$.

	For the case of $\tau^{\geq n}$, instead define
	$\underline{h}^{n+1} := \bigl(Y^n \twoheadrightarrow
	\Coker(d_Y^{n-1})\bigr) \circ h^{n+1}$ and use this to define
	$\underline{h} \in
	\Hom^{-1}\left( \tau^{\geq n} X, \tau^{\geq n} Y \right)$; the
	remainder is similar. The two cases are of course dual.
\end{proof}

% segment-id: o014.li2.chapter3.truncation-functors.p011
The following simple but important property follows immediately from
Definition~\ref{def:truncation-cplx}.

% segment-id: o014.li2.chapter3.truncation-functors.pr002
\begin{proposition}\label{prop:truncate-twice}
	For all integers $a<b$ and $n$, and
	$X \in \Obj(\cate{C}(\mathcal{A}))$, we have\footnote{Translator's note
	(O014-C053): the source writes $X \in \Obj(\mathcal{A})$, but the
	truncation functors and $\Hm^n$ in this proposition act on complexes.}
% segment-id: o014.li2.chapter3.truncation-functors.q008
	\begin{gather*}
		\tau^{\leq a} \tau^{\geq b}(X) = 0 = \tau^{\geq b} \tau^{\leq a}(X), \\
		\tau^{\leq n }\tau^{\geq n}(X) \simeq \Hm^n(X)[-n] \simeq \tau^{\geq n} \tau^{\leq n}(X) \quad \text{(canonical isomorphisms).}
	\end{gather*}
\end{proposition}
